\documentclass[11pt]{article}
\usepackage{graphicx}
\usepackage[bookmarks=true]{hyperref}
\usepackage{bookmark}
\usepackage{hyperref}
\usepackage{csquotes}
\usepackage{float}
\usepackage{wrapfig}
\usepackage[normalem]{ulem}
\useunder{\uline}{\ul}{}

\usepackage{array}
\newcolumntype{L}[1]{>{\raggedright\let\newline\\\arraybackslash\hspace{0pt}}m{#1}}
\newcolumntype{C}[1]{>{\centering\let\newline\\\arraybackslash\hspace{0pt}}m{#1}}
\newcolumntype{R}[1]{>{\raggedleft\let\newline\\\arraybackslash\hspace{0pt}}m{#1}}

\setlength{\parindent}{0pt}

\begin{document}

\begin{titlepage}
\begin{flushright}

\includegraphics[width=380px]{../global/University_of_Pretoria_Logo.png}
\newline
\newline

\textbf {\LARGE Plan for Software Aspects of Certification} \newline

\textbf {\Large (PSAC)} \newline

\centering\includegraphics[width=100px]{../global/Logo.jpg}

\textbf {\Large Linphone for Android Group Chat (Waterfall)}\newline

\flushright \textbf {\large Version: 1.2}\newline

\centering \textbf {\large Authors:}

\begin{table}[H]
\large
\centering
\begin{tabular}{rl}
	Izak Blom & 13126777 \\
	David Breetzke & 12056503 \\
	Paul Engelke & 13093500 \\
	Prenolan Govender & 13102380 \\
	Jessica Lessev & 13049136 \\
\end{tabular}
\end{table}

Date: \today

\end{flushright}
\end{titlepage}

\setcounter{tocdepth}{3}
\setcounter{secnumdepth}{5}
\tableofcontents

\newpage
\section{Revision History}
\begin{table}[h]
\begin{tabular}{llll}
\textbf{Date}          & \textbf{Description}  & \textbf{Author}       & \textbf{Comments}   \\ \hline
\multicolumn{1}{|R{2cm}|}{26/08/2015} & \multicolumn{1}{L{4.5cm}|}{Document Creation} & \multicolumn{1}{l|}{Team Eclectic} & \multicolumn{1}{L{4cm}|}{Version 1} \\ \hline
\multicolumn{1}{|l|}{} & \multicolumn{1}{l|}{} & \multicolumn{1}{l|}{} & \multicolumn{1}{l|}{} \\ \hline
\multicolumn{1}{|l|}{} & \multicolumn{1}{l|}{} & \multicolumn{1}{l|}{} & \multicolumn{1}{l|}{} \\ \hline
\multicolumn{1}{|l|}{} & \multicolumn{1}{l|}{} & \multicolumn{1}{l|}{} & \multicolumn{1}{l|}{} \\ \hline
\multicolumn{1}{|l|}{} & \multicolumn{1}{l|}{} & \multicolumn{1}{l|}{} & \multicolumn{1}{l|}{} \\ \hline
\multicolumn{1}{|l|}{} & \multicolumn{1}{l|}{} & \multicolumn{1}{l|}{} & \multicolumn{1}{l|}{} \\ \hline
\multicolumn{1}{|l|}{} & \multicolumn{1}{l|}{} & \multicolumn{1}{l|}{} & \multicolumn{1}{l|}{} \\ \hline
\multicolumn{1}{|l|}{} & \multicolumn{1}{l|}{} & \multicolumn{1}{l|}{} & \multicolumn{1}{l|}{} \\ \hline
\multicolumn{1}{|l|}{} & \multicolumn{1}{l|}{} & \multicolumn{1}{l|}{} & \multicolumn{1}{l|}{} \\ \hline
\multicolumn{1}{|l|}{} & \multicolumn{1}{l|}{} & \multicolumn{1}{l|}{} & \multicolumn{1}{l|}{} \\ \hline
\multicolumn{1}{|l|}{} & \multicolumn{1}{l|}{} & \multicolumn{1}{l|}{} & \multicolumn{1}{l|}{} \\ \hline
\multicolumn{1}{|l|}{} & \multicolumn{1}{l|}{} & \multicolumn{1}{l|}{} & \multicolumn{1}{l|}{} \\ \hline
\multicolumn{1}{|l|}{} & \multicolumn{1}{l|}{} & \multicolumn{1}{l|}{} & \multicolumn{1}{l|}{} \\ \hline
\multicolumn{1}{|l|}{} & \multicolumn{1}{l|}{} & \multicolumn{1}{l|}{} & \multicolumn{1}{l|}{} \\ \hline
\multicolumn{1}{|l|}{} & \multicolumn{1}{l|}{} & \multicolumn{1}{l|}{} & \multicolumn{1}{l|}{} \\ \hline
\end{tabular}
\end{table}

\section{Document Approval}
\begin{table}[h]
\begin{tabular}{llll}
\textbf{Signature}     & \textbf{Printed Name} & \textbf{Title}        & \textbf{Comments}     \\ \hline
\multicolumn{1}{|l|}{} & \multicolumn{1}{L{3.5cm}|}{} & \multicolumn{1}{L{3.5cm}|}{} & \multicolumn{1}{L{4cm}|}{} \\ \hline
\multicolumn{1}{|l|}{} & \multicolumn{1}{l|}{} & \multicolumn{1}{l|}{} & \multicolumn{1}{l|}{} \\ \hline
\multicolumn{1}{|l|}{} & \multicolumn{1}{l|}{} & \multicolumn{1}{l|}{} & \multicolumn{1}{l|}{} \\ \hline
\multicolumn{1}{|l|}{} & \multicolumn{1}{l|}{} & \multicolumn{1}{l|}{} & \multicolumn{1}{l|}{} \\ \hline
\end{tabular}
\end{table}

\newpage

\section{Introduction}
\subsection{Purpose}
This document describes the plan for testing the development and functionality of the extended Linphone application software against the
user requirements as defined in the software requirements specification document [SRS]. The purpose of this acceptance test is to make sure
that the system developed during the Linphone extension project complies with the requirements listed in the SRS. These tests are executed in the Acceptance Test (AT) phase of the project.

\subsection{Scope}
The Linphone application is based on a distributed system. It ranges from a wide variety of technologies and incorporates a wide range of aspects. The extension to Linphone enables users to create, manage and participate in group chats with the side choice of using encryption while doing so. By hiding all complexity from the main server it makes the system easier to use, implement and test. Usability is also increased by providing an interactive user interface as a front-end for the users to interact with.

\subsection{Acronyms \& Abbreviations}
\begin{itemize}
\item \textbf{ATP} Acceptance Test Procedure
\item \textbf{AT} Acceptance Test
\item \textbf{SRS} Software Requirements Specification
\item \textbf{SDD} Software Design Description
\item \textbf{AES256} Advanced Encryption Standard (256-bit)
\end{itemize}
\subsection{Applicable Documents}
\subsubsection{Internal Documents}
\begin{tabular}{ll}
\textbf{}       & \textbf{}  \\ \hline
\multicolumn{1}{|L{2cm}|}{SRS} & \multicolumn{1}{L{10cm}|}{Software Requirements Specification} \\ \hline
\multicolumn{1}{|L{2cm}|}{PSAC} & \multicolumn{1}{L{10cm}|}{Plan for Software Aspects of Certification} \\ \hline
\multicolumn{1}{|L{2cm}|}{SDD} & \multicolumn{1}{L{10cm}|}{Software Design Description} \\ \hline
\end{tabular}

\subsection{Test Environment}
Several aspects play a role in the testing environment. These include equipment used to test the product, technologies used, the environmental/surrounding conditions etc...\\
\textbf{Equipment}\\
The equipment used to test the extended functionality for the project will be as follows:
\begin{itemize}
\item Computer/Laptop with necessary IDE's and testing environment
\item Provided cellular phones with appropriate operating system and application installation
\end{itemize}
\textbf{Technologies}\\
The technologies used to test the functionality are mainly software oriented and include:
\begin{itemize}
\item JUnit for unit testing and integration testing
\item .............................
\end{itemize}

\textbf{Test Conditions}
Several factors describe the test conditions:
\begin{itemize}
\item Testing will take place in an integrated environment
\item The project is developed in a Linux environment using the Eclipse IDE. Thus testing will take place within the development paradigm using software to assist in unit and integration testing.
\item Since the project is a development for an android device, the bulk of the testing will occur on an android device. The specific device to be used will be the Zest cell phone.
\item All necessary actions are performed. In order to see the application in a certain state, we need to actually perform every step to get us there (clicking buttons, filling in text boxes, sending messages etc).
\item The only thing that is mocked during the testing is data such as message content etc
\end{itemize}



\subsection{Equipment Configuration \& Test Case Setup Instructions}


\section{Test Procedures \& Verification Results Record}
\subsection{Test Procedures Summary}
All tests will involve testing the use case preconditions and postconditions.
The following is a list of tests are to be performed:
\begin{enumerate}
\item Create group chat
\item Send Invite
\item Delete Group Chat
\item Update Group Chat
\begin{enumerate}
\item Update group description
\item Add Member
\item Remove Member
\item Change Administrator 
\end{enumerate}
\item Send Message
\item Receive Message
\end{enumerate}



\subsection{Test 1: Create group chat}
\subsubsection{Description}
A user shall be able to create a new group.\\
\textbf{Preconditions:} For each of the pre-conditions below an exception is raised where that precondition is not met, to indicate service refusal.
The group will not be created if one of the following preconditions is not met:
\begin{itemize}
\item User must have access to the Linphone service.
\item The user may not create the group unless at least one other member is added.
\end{itemize}
\textbf{Postconditions:} The post-conditions specify the conditions which must hold true when the service has been provided.
\begin{itemize}
\item The creator is automatically assigned administrator rights.
\item An invite is sent to the other added users to ask them if they want to be part of the group.
\item The group is created with the respective profiles and members.
\item An error is displayed if the user attempts to create a group with less than two members in it.
\end{itemize}

\subsubsection{Expected Result to Pass Test}
In order for the application to pass this test a new group chat should be created without any unexpected problems. The appropriate preconditions and postconditions should be followed.
\subsubsection{Result}


\subsection{Test 2: Send Invite}
An invite shall be sent to each user added to the group upon creation or when a new member is added to the group. Thus this test is divided into two sub-tests.

\subsubsection{Test 2.1: Invite sent on group chat creation}
\paragraph{Description}
 A new group chat will be created and checked whether the initially added participants receive the invite to join the group chat.\\
 \textbf{Preconditions:} For each of the pre-conditions below an exception is raised where that precondition is not met, to indicate service refusal.
 The invite will not be sent if one of the following scenarios occurs:
 \begin{itemize}
 \item The receiver of the invite is a contact added by the administrator of the group.
 \end{itemize}
 \textbf{Postconditions:} The post-conditions specify the conditions which must hold true when the service has been provided.
 \begin{itemize}
 \item The invite must be sent to the correct user.
 \end{itemize}
\paragraph{Expected Result to Pass Test}
In order for the application to pass this test the members added to the new group chat created must receive an invite to participate.
\paragraph{Result}

\subsubsection{Test 2.2: Invite sent on group chat addition of member}
\paragraph{Description}
 An existing group chat will be edited by adding a new member to participate in the group chat. It will be checked whether the newly added member receives an invite for the group chat. \\
  \textbf{Preconditions:} For each of the pre-conditions below an exception is raised where that precondition is not met, to indicate service refusal.
  The invite will not be sent if one of the following scenarios occurs:
  \begin{itemize}
  \item The receiver of the invite is a contact added by the administrator of the group.
  \end{itemize}
  \textbf{Postconditions:} The post-conditions specify the conditions which must hold true when the service has been provided.
  \begin{itemize}
  \item The invite must be sent to the correct user.
  \end{itemize}
\paragraph{Expected Result to Pass Test}
In order for the application to pass this test the member added after the group chat was created must receive an invite to participate.
\paragraph{Result}

\subsection{Test 3: Delete Group Chat}
\subsubsection{Description}
A user shall be able to leave a group chat. When all the users leave a group chat, it will be deleted. If only one person remains in a group chat, that person will be unable to send messages to the chat until at least one other person is added to the chat. If the group owner leaves the group, a new owner is automatically assigned.\\
\textbf{Preconditions:} 
\begin{itemize}
\item The user must be part of a group.
\end{itemize}
\textbf{Postconditions:}
\begin{itemize}
\item The chat is locally deleted from the user's phone.
\item The user is removed/unlinked from the group chat.
\item If the user leaving is the administrator: The administrative privileges are assigned to another group member.
\item If the user is the last user in the group, their ability to send messages  to the group is locked.
\end{itemize}
\subsubsection{Expected Result to Pass Test}
In order for the application to pass this test the group chat in question must be deleted and removed from the members profile. All post-conditions must take place.
\subsubsection{Result}

\subsection{Test 4: Update Group Chat}
The group name, subject, and members may be changed. Thus this use case is divided into three major sub-tasks (Update group chat descriptions - name, subject etc, add member, remove member) with minor subtask within each of the major ones.
\subsubsection{Test 4.1: Update Group descriptions}
\paragraph{Description} Here a group description such as a name or subject will be changed and saved. It follows that each member on the corresponding group chat's profile will be updated with the correct information.
\paragraph{Expected Result to Pass Test}
Descriptions updated successfully in all relevant places.\\
\textbf{Preconditions:} 
\begin{itemize}
\item A group exists and has members.
\end{itemize}
\textbf{Postconditions:}
\begin{itemize}
\item The current name or subject has been changed.
\end{itemize}
\paragraph{Result}

\subsubsection{Test 4.2: Add Member}
\paragraph{Description}
The administrator shall be able to add new users to the group. Once a new member is added it will be tested that the member receives the invite and is linked to the specific group chat.\\
\textbf{Preconditions:} 
\begin{itemize}
\item Group exists with an administrator in it.
\end{itemize}
\textbf{Postconditions:}
\begin{itemize}
\item The user is added to the list of current users in the group and receives messages, etc.
\end{itemize}
\paragraph{Expected Result to Pass Test}
New member receives and invite to participate in group chat, is linked to the group chat, can send and receive messages on the group.
\paragraph{Result}
\subsubsection{Test 4.2: Remove Member}
\paragraph{Description}
The administrator shall be able to remove members from the group. That user is removed from the list of users currently in the group.\\
\textbf{Preconditions:} 
\begin{itemize}
\item A group exists with more than one member.
\end{itemize}
\textbf{Postconditions:}
\begin{itemize}
\item The user is no longer a part of the group which entails the user not receiving messages from the group, etc.
\end{itemize}
\paragraph{Expected Result to Pass Test}
The user no longer has access to the group chat, thus cannot participate. All left over members' member list within the group chat is updated to one less member.
\paragraph{Result}

\subsubsection{Test 4.3: Change Administrator}
\paragraph{Description}
The administrator shall be able to reassign a group
administrator. This removes administrator privileges from the current administrator and reassigns them to another user within the group.\\
\textbf{Preconditions:} 
\begin{itemize}
\item Group exists with an administrator and at least
one other user.
\end{itemize}
\textbf{Postconditions:}
\begin{itemize}
\item The user selected receives administrator privileges while the current administrator is stripped of the privileges.
\end{itemize}
\paragraph{Expected Result to Pass Test}
The previous administrator no longer had administrative rights, the new administrator has administrative rights, all group chat profiles are updated to show the new administrator.
\paragraph{Result}

\subsection{Test 5: Send Message}
\subsubsection{Description}
A user shall be able to send a message to the group.
The message will be broadcast to all members of the group.\\
\textbf{Preconditions:} 
\begin{itemize}
\item The user is a member of the group chat.
\item The interface that has focus, is the group chat interface.
\end{itemize}
\textbf{Postconditions:}
\begin{itemize}
\item The user's message has been sent to all other participants.
\end{itemize}
\subsubsection{Expected Result to Pass Test}
The message appears as sent at the senders end and is received by the rest of the group members.
\subsubsection{Result}

\subsection{Test 6: Receive Message}
\subsubsection{Description}
A message object sent on a group shall be received
by every member subscribed to a group.\\
\textbf{Preconditions:} 
\begin{itemize}
\item A message object has been broadcast to all users subscribed to a particular group.
\item The user receiving an incoming message should be subscribed to the
group from which the message is broadcast.
\end{itemize}
\textbf{Postconditions:}
\begin{itemize}
\item The message will be displayed on the group chat interface
\end{itemize}
\subsubsection{Expected Result to Pass Test}
All sent messages are received on the corresponding group chat members devices.
\subsubsection{Result}

\section{Summary Test Report}
\begin{table}[h]
\begin{tabular}{ll}
\textbf{Test}     & \textbf{Pass/Fail} \\ \hline
\multicolumn{1}{|L{10cm}|}{Test 1} & \multicolumn{1}{L{2cm}|}{} \\ \hline
\multicolumn{1}{|L{10cm}|}{Test 2} & \multicolumn{1}{L{2cm}|}{} \\ \hline
\multicolumn{1}{|L{10cm}|}{Test 3} & \multicolumn{1}{L{2cm}|}{} \\ \hline
\multicolumn{1}{|L{10cm}|}{Test 4} & \multicolumn{1}{L{2cm}|}{} \\ \hline
\multicolumn{1}{|L{10cm}|}{Test 4.1} & \multicolumn{1}{L{2cm}|}{} \\ \hline
\multicolumn{1}{|L{10cm}|}{Test 4.2} & \multicolumn{1}{L{2cm}|}{} \\ \hline
\multicolumn{1}{|L{10cm}|}{Test 4.3} & \multicolumn{1}{L{2cm}|}{} \\ \hline
\multicolumn{1}{|L{10cm}|}{Test 4.4} & \multicolumn{1}{L{2cm}|}{} \\ \hline
\multicolumn{1}{|L{10cm}|}{Test 5} & \multicolumn{1}{L{2cm}|}{} \\ \hline
\multicolumn{1}{|L{10cm}|}{Test 6} & \multicolumn{1}{L{2cm}|}{} \\ \hline
\end{tabular}
\end{table}

\section{Traceability Matrix}
\end{document}